\begin{abstract}
Standard-cell timing characterization remains a major bottleneck in advanced-node VLSI design due to the high cost of transistor-level simulation and the scarcity of labeled data at target technology nodes. This paper presents a cross-technology-node timing prediction framework for standard-cell libraries that combines heterogeneous graph learning with transfer learning. We represent each cell as a heterogeneous circuit graph with explicit transistor, resistor, capacitor, and port nodes, enabling joint modeling of topology and parasitics. Based on this representation, we propose an H-GAT encoder with structure-guided attention and physics-aware message passing to capture both logical interactions and physical effects. To alleviate domain shift across process nodes, we further disentangle circuit representations into technology-invariant structural features and technology-dependent physical features, and introduce a transfer strategy that preserves shared structural knowledge while aligning process-related features across domains. Experiments on cross-node transfer from Nangate45 to ASAP7 show that our method outperforms target-only training, direct source-target mixing, and pre-training plus fine-tuning baselines under low-resource target-node settings. In particular, our approach achieves [XX]\% lower MAE / [XX]\% lower MAPE, improves $R^2$ to [XX], and delivers [XX]$\times$ speedup over HSPICE-based characterization, demonstrating its effectiveness for fast and accurate advanced-node timing prediction.
\end{abstract}

